% ============================================================
%  CHƯƠNG 5: SO SÁNH GMM VỚI K-MEANS
%  Người viết: [Họ tên thành viên phụ trách]
% ============================================================

\section{So sánh GMM với \texorpdfstring{$k$}{k}-means}
\label{sec:kmeans_vs_gmm}

% -----------------------------------------------------------
\subsection{Ôn lại thuật toán \texorpdfstring{$k$}{k}-means}
% -----------------------------------------------------------
% Gợi ý nội dung:
%   - Mô tả bước gán nhãn (assignment) và cập nhật tâm cụm (update)
%   - Hàm mục tiêu: minimize Σ_n Σ_k r_{nk} ||x_n - μ_k||²
%   - Đây là phân cụm cứng (hard assignment): r_{nk} ∈ {0,1}

% TODO: Viết nội dung ở đây


% -----------------------------------------------------------
\subsection{\texorpdfstring{$k$}{k}-means là trường hợp giới hạn của GMM-EM}
% -----------------------------------------------------------
% Gợi ý nội dung:
%   - Khi Σ_k = ε²I và ε → 0, responsibility γ_{nk} → r_{nk} ∈ {0,1}
%   - M-step GMM trong giới hạn này trùng với bước cập nhật tâm k-means
%   - Phân tích toán học: dẫn xuất từ công thức softmax của γ_{nk}

% TODO: Viết nội dung ở đây


% -----------------------------------------------------------
\subsection{Phân cụm mềm vs. phân cụm cứng}
% -----------------------------------------------------------
% Gợi ý nội dung:
%   - Ưu điểm phân cụm mềm (GMM): uncertainty, điểm nằm giữa cụm
%   - Hạn chế phân cụm cứng (k-means): mất thông tin biên giới cụm
%   - Ví dụ minh họa bằng hình vẽ hoặc dữ liệu 2D

% TODO: Viết nội dung ở đây

% Ví dụ chèn hình so sánh:
% \begin{figure}[H]
%   \centering
%   \includegraphics[width=0.85\linewidth]{figures/hard_vs_soft.pdf}
%   \caption{So sánh phân cụm cứng ($k$-means) và phân cụm mềm (GMM)}
%   \label{fig:hard_vs_soft}
% \end{figure}


% -----------------------------------------------------------
\subsection{Bảng tổng hợp so sánh}
% -----------------------------------------------------------
% Gợi ý nội dung: điền bảng so sánh các tiêu chí sau

\begin{table}[H]
  \centering
  \caption{So sánh $k$-means và GMM-EM}
  \label{tab:kmeans_vs_gmm}
  \begin{tabularx}{\linewidth}{lXX}
    \toprule
    \textbf{Tiêu chí}         & \textbf{$k$-means}          & \textbf{GMM-EM} \\
    \midrule
    Loại phân cụm             & Cứng (hard)                 & Mềm (soft) \\
    Hàm mục tiêu              & Tổng bình phương khoảng cách & Log-likelihood \\
    Mô hình xác suất          & Không                       & Có \\
    Hình dạng cụm             & Hypersphere (isotropic)     & Ellipsoid tùy ý \\
    Độ phức tạp mỗi bước      & $O(NKD)$                    & $O(NKD^2 + KD^3)$ \\
    Thông tin bất định        & Không                       & Có (qua $\gamma_{nk}$) \\
    Nhạy với khởi tạo         & Cao                         & Cao \\
    % TODO: bổ sung thêm hàng nếu cần
    \bottomrule
  \end{tabularx}
\end{table}

% TODO: Thêm nhận xét/phân tích bảng trên


% -----------------------------------------------------------
\subsection{Khi nào dùng GMM, khi nào dùng \texorpdfstring{$k$}{k}-means?}
% -----------------------------------------------------------
% Gợi ý nội dung:
%   - Dữ liệu có cụm hình ellipsoid → GMM
%   - Cần xác suất thuộc cụm (uncertainty) → GMM
%   - Tập dữ liệu lớn, cần tốc độ → k-means (hoặc mini-batch k-means)
%   - Số chiều rất cao → cân nhắc cả hai, regularize Σ_k

% TODO: Viết nội dung ở đây
